\clearpage
\subsection*{Part (b)}
\addcontentsline{toc}{subsection}{Problem 6(b)}
\markright{Problem 6: Part (b)}

\textbf{Problem Statement (Textbook 4.2.1b):} 

Prove these facts, needed in the proof of Theorem 2.

\textbf{Problem 6(b)} The inverse of a unit lower triangular matrix is unit lower triangular.

\noindent\rule{\textwidth}{0.4pt}

\vspace{1em}

\textbf{Solution:}

Let $L$ be unit lower triangular (lower triangular with ones on the diagonal). To find $L^{-1}$, we solve $LX = I$ column by column. For the $j$-th column $\mathbf{x}_j$ of $L^{-1}$, we solve:
$$L\mathbf{x}_j = \mathbf{e}_j$$

Since $L$ is unit lower triangular, this system can be solved by forward substitution starting from row 1. The key observation is:
\begin{itemize}
    \item Row 1: $x_{j1} = \delta_{j1}$ (so $x_{j1} = 1$ if $j=1$, else $x_{j1} = 0$)
    \item Row 2: $l_{21}x_{j1} + x_{j2} = \delta_{j2}$ gives $x_{j2} = 0$ for all $j > 2$
    \item Continuing downward, row $i$ depends only on $x_{j1}, x_{j2}, \ldots, x_{ji}$
\end{itemize}

By induction, for $i < j$, we have $x_{ji} = 0$ (zeros above diagonal). Moreover, since row $j$ gives $l_{j1}x_{j1} + \cdots + l_{j,j-1}x_{j,j-1} + x_{jj} = 1$, we get $x_{jj} = 1$ (ones on diagonal). Therefore, $L^{-1}$ is unit lower triangular. $\square$

\vspace{1em}

\textbf{Example:} Consider $L = \begin{bmatrix} 1 & 0 & 0 \\ 2 & 1 & 0 \\ 1 & 3 & 1 \end{bmatrix}$. To find $L^{-1}$, solve $L\mathbf{x}_j = \mathbf{e}_j$ for each column:

\textit{Column 1:} $L\mathbf{x}_1 = \begin{bmatrix} 1 \\ 0 \\ 0 \end{bmatrix}$
\begin{align*}
\text{Row 1:} \quad 1 \cdot x_{11} &= 1 \implies x_{11} = 1 \\
\text{Row 2:} \quad 2(1) + 1 \cdot x_{21} &= 0 \implies x_{21} = -2 \\
\text{Row 3:} \quad 1(1) + 3(-2) + 1 \cdot x_{31} &= 0 \implies x_{31} = 5
\end{align*}
Thus $\mathbf{x}_1 = [1, -2, 5]^T$ (ones on diagonal, zeros above).

\textit{Column 2:} $L\mathbf{x}_2 = \begin{bmatrix} 0 \\ 1 \\ 0 \end{bmatrix}$
\begin{align*}
\text{Row 1:} \quad x_{12} &= 0 \\
\text{Row 2:} \quad 2(0) + x_{22} &= 1 \implies x_{22} = 1 \\
\text{Row 3:} \quad 1(0) + 3(1) + x_{32} &= 0 \implies x_{32} = -3
\end{align*}
Thus $\mathbf{x}_2 = [0, 1, -3]^T$ (zero in position $(1,2)$, one on diagonal).

\textit{Column 3:} $L\mathbf{x}_3 = \begin{bmatrix} 0 \\ 0 \\ 1 \end{bmatrix}$
\begin{align*}
\text{Row 1:} \quad x_{13} &= 0 \\
\text{Row 2:} \quad 2(0) + x_{23} &= 0 \implies x_{23} = 0 \\
\text{Row 3:} \quad 1(0) + 3(0) + x_{33} &= 1 \implies x_{33} = 1
\end{align*}
Thus $\mathbf{x}_3 = [0, 0, 1]^T$ (zeros above diagonal, one on diagonal).

Therefore: $L^{-1} = \begin{bmatrix} 1 & 0 & 0 \\ -2 & 1 & 0 \\ 5 & -3 & 1 \end{bmatrix}$ is unit lower triangular.

\vspace{2em}
\noindent\rule{\textwidth}{0.4pt}
